% ============================================================
%  MMC e MDC — Apostila Premium | PratiqueJá
%  Engine: XeLaTeX
% ============================================================
\documentclass[12pt]{article}
\usepackage{../../pratiqueja}

% \hlm — destaque de resultado em modo-math (cor sem bold)
\newcommand{\hlm}[1]{{\color{darkblue}#1}}

% Tabela de fatoração simultânea (a "chave")
% Uso: \begin{fatora} <linhas> \end{fatora}
\newenvironment{fatora}{%
  \par\vspace{2pt}\begin{center}
  \renewcommand{\arraystretch}{1.25}%
  \small
}{%
  \end{center}\vspace{2pt}\par
}

\setsubject{MMC e MDC}

\begin{document}
\pagenumbering{arabic}
\setstretch{1.2}
\color{bodytext}

\heroheader{MMC e MDC}%
           {Mínimo Múltiplo e Máximo Divisor Comum}%
           {Matemática · Básico}

\setlength{\columnsep}{18pt}
\setlength{\columnseprule}{0.5pt}
\def\columnseprulecolor{\color{exborder}}

\begin{multicols}{2}

%─────────────────────────────────────────────────────────────
\TW{MMC}{badgepurple}{Mínimo Múltiplo Comum}{Definição e método de cálculo}

\regra{O \textbf{MMC} de dois ou mais inteiros positivos é o
\textbf{menor inteiro diferente de zero} que é múltiplo de todos
eles ao mesmo tempo.}

\small\color{bodytext}%
$a$ é múltiplo de $b$ quando existe inteiro $q$ com $a = b\cdot q$.
Ex.: $10 = 5\cdot2$ e $18 = 6\cdot3$.

\regra{\textbf{Fatoração simultânea:} divida pelo menor primo que
divida \emph{ao menos um} dos números, carregando os demais
inalterados. Repita até todos chegarem a $1$. O MMC é o
\textbf{produto de todos os fatores primos usados}.}

\obs{Pré-requisito: decomposição em fatores primos — ver apostila
\textit{Números Primos}.}

%─────────────────────────────────────────────────────────────
\TW{2}{darkblue}{MMC com 2 números}{Exemplo — MMC de 12 e 18}

Quando um número não é divisível pelo fator, ele é \textbf{carregado}
para a linha seguinte sem alteração.

\begin{fatora}
\begin{tabular}{r@{,\;}l|l}
12 & 18 & 2\\
 6 &  9 & 2 \;{\footnotesize\color{mutedtext}(9 carregado)}\\
 3 &  9 & 3\\
 1 &  3 & 3 \;{\footnotesize\color{mutedtext}(1 carregado)}\\
\hline
 1 &  1 & \\
\end{tabular}
\end{fatora}

\exemp{MMC$(12,18) = 2^2\cdot3^2 = 4\cdot9 = \hlm{36}$}

%─────────────────────────────────────────────────────────────
\TW{3}{darkblue}{MMC com 3 números}{Exemplo — MMC de 90, 30 e 25}

\begin{fatora}
\begin{tabular}{r@{,\;}r@{,\;}l|l}
90 & 30 & 25 & 2\\
45 & 15 & 25 & 3\\
15 &  5 & 25 & 3\\
 5 &  5 & 25 & 5\\
 1 &  1 &  5 & 5\\
\hline
 1 &  1 &  1 & \\
\end{tabular}
\end{fatora}

\exemp{MMC$(90,30,25) = 2\cdot3^2\cdot5^2 = 2\cdot9\cdot25 = \hlm{450}$}

%─────────────────────────────────────────────────────────────
\TW{MDC}{badgepurple}{Máximo Divisor Comum}{Definição e método de cálculo}

\regra{O \textbf{MDC} de dois ou mais inteiros é o \textbf{maior
inteiro positivo} que divide todos eles sem deixar resto.}

\regra{\textbf{Decomposição simultânea:} divida pelo menor primo que
divida \emph{todos} ao mesmo tempo. \textbf{Pare} quando não existir
mais primo que divida todos juntos. O MDC é o produto dos fatores
usados.}

\obs{Diferença para o MMC: no MDC o fator só vale se dividir
\emph{todos} os números da linha. As linhas \hl{destacadas} abaixo
são essas.}

%─────────────────────────────────────────────────────────────
\TW{2}{badgegreen}{MDC com 2 números}{Exemplo — MDC de 32 e 24}

Dividimos enquanto o fator divide \textbf{os dois}. Quando isso deixa
de ser possível, paramos.

\begin{fatora}
\begin{tabular}{r@{,\;}l|l}
\rowcolor{rulebg}32 & 24 & 2\\
\rowcolor{rulebg}16 & 12 & 2\\
\rowcolor{rulebg} 8 &  6 & 2\\
\hline
 4 &  3 & \\
\end{tabular}
\end{fatora}

$4$ e $3$ não têm fator primo comum — paramos aqui.

\exemp{MDC$(32,24) = 2^3 = \hlm{8}$}

%─────────────────────────────────────────────────────────────
\TW{3}{badgegreen}{MDC com 3 números}{Exemplo — MDC de 120, 30 e 45}

Como $45$ é ímpar, nenhuma divisão por $2$ serve aos três —
começamos pelo $3$.

\begin{fatora}
\begin{tabular}{r@{,\;}r@{,\;}l|l}
\rowcolor{rulebg}120 & 30 & 45 & 3\\
\rowcolor{rulebg} 40 & 10 & 15 & 5\\
\hline
  8 &  2 &  3 & \\
\end{tabular}
\end{fatora}

$8$, $2$ e $3$ não têm fator comum — paramos.

\exemp{MDC$(120,30,45) = 3\cdot5 = \hlm{15}$}

%─────────────────────────────────────────────────────────────
\TW{$\star$}{darkblue}{Propriedades}{Relações entre MMC e MDC}

\formula{MMC$(a,b)\times$MDC$(a,b) = a\times b$}

\SH{Verificação com 12 e 18}
\exemp{%
  MDC$(12,18) = \hlm{6}$ \;{\footnotesize(comuns: $2\cdot3$)}\\
  MMC$(12,18) = \hlm{36}$\\
  $36\times6 = 216 = 12\times18$ \ok%
}

\SH{Caso especial: divisibilidade direta}
\exemp{%
  \textbf{Se $a$ divide $b$:} MDC$(a,b)=a$ e MMC$(a,b)=b$\\[2pt]
  Ex.: MDC$(5,30)=\hlm{5}$ e MMC$(5,30)=\hlm{30}$%
}

%─────────────────────────────────────────────────────────────
\TW{Ap}{badgegreen}{Aplicação Prática}{Quando usar MMC e quando usar MDC}

\regra{Use o \textbf{MMC} para o \emph{menor momento em que dois
ciclos coincidem} (encontros, sinais, repetições).}

\SH{Problema — encontro de ônibus}
\exemp{%
  Um ônibus passa a cada $12$\,min e outro a cada $18$\,min,
  partindo juntos. Quando voltam a partir juntos?\\[2pt]
  MMC$(12,18) = 2^2\cdot3^2 = \hlm{36\text{ minutos}}$%
}

\regra{Use o \textbf{MDC} para o \emph{maior tamanho ou quantidade
igual} que cabe exatamente em dois ou mais totais (sem sobras).}

\SH{Problema — divisão em grupos}
\exemp{%
  Dividir $32$ meninos e $24$ meninas em grupos de mesmo tamanho,
  sem misturar. Maior grupo possível?\\[2pt]
  MDC$(32,24) = 2^3 = \hlm{8\text{ alunos por grupo}}$\\[2pt]
  {\footnotesize Resultado: 4 grupos de meninos + 3 de meninas.}%
}

\end{multicols}

\end{document}
